% 28 main slides. Optional explanations: slides 6 and 14.

\begin{frame}[plain]
\vspace{.7cm}
{\color{navy}\fontsize{25}{31}\selectfont\bfseries
When Will a Borrower Default?\par}
\vspace{.45cm}
{\Large Survival Analysis in Credit Risk\par}
\vspace{1.2cm}
{\large\color{teal}Luke Munnelly}
\talknote{0:20}{Introduce the lending question rather than the method. This is a standalone lecture for MSc Stochastics and Financial Mathematics students. We will use a small audited example to interpret default rates, not compare prediction algorithms. Advance promptly to the two borrowers.}
\end{frame}

\begin{frame}{Should these be the same outcome?}
\begin{itemize}\setlength\itemsep{.65cm}
\item<1-> One borrower defaults after three months.
\item<2-> Another defaults after two years.
\end{itemize}
\vspace{.6cm}
\uncover<3->{\textcolor{teal}{Both defaulted. The timing is different.}}\par
\vspace{.4cm}\qual{Hypothetical loans.}
\talknote{0:55}{Start with the exact question: One borrower defaults after three months, another after two years. Should we treat these as the same outcome? Show only the first borrower, then reveal the second and pause for three seconds. They share an eventual default label but expose the lender for very different lengths of time. Cash flows and time exposed to credit risk differ; no numerical loss calculation is needed. Reveal the final statement only after making this distinction. This is a rhetorical pause, not planned audience-question time.}
\end{frame}

\begin{frame}{Principal still outstanding}
\qual{Illustrative: \texteuro12,000 loan; \texteuro400 principal repaid each month.}
\vspace{.55cm}
\begin{center}
\begin{tabular}{lr}
After 3 payments & \uncover<2->{\texteuro10,800 outstanding}\\[.45cm]
After 24 payments & \uncover<2->{\texteuro2,400 outstanding}
\end{tabular}
\end{center}
\vspace{.5cm}
\uncover<3->{Earlier default can occur before much principal is repaid.}\par
\vspace{.3cm}\qual{Assume default just after payment; ignore interest, fees and recoveries.}
\talknote{0:50}{Use a deliberately hypothetical principal-only schedule: twelve thousand euros, four hundred euros principal repaid each month. Explain the setup before revealing both balances. After three payments, 12000 minus 3 times 400 leaves 10800; after twenty-four payments, 2400 remains. Reveal the final interpretation only after the arithmetic. Earlier default can leave more principal outstanding, which makes timing economically relevant. This is not a calculation of actual credit loss: recoveries, interest and other factors are deliberately omitted. The hypothetical euro units are chosen for teaching, not inferred from Bondora.}
\end{frame}

\begin{frame}{What do six months tell us?}
\centering
Observed for six months; no default so far\par
\vspace{.65cm}
\uncover<2->{{\Large $T_D>6$ months}\par
\vspace{.45cm}Contributes to risk sets while observed.}
\talknote{0:50}{For this actively observed hypothetical loan, six default-free months tell us T_D is greater than six months, not what its eventual value will be. Reveal the inequality and the risk-set contribution together after describing the history. This loan contributes to the exposed risk sets until its observation ends; it is not discarded and is not assumed to remain default-free forever. A closure would instead remove its exposure through a competing event. We cannot assign a certain individual default date from this partial history.}
\end{frame}

\begin{frame}{A binary label needs a horizon}
\begin{itemize}\setlength\itemsep{.6cm}
\item<1-> A classifier can predict default within 12 months.
\item<2-> Six months without default does not establish a 12-month label.
\end{itemize}
\vspace{.5cm}
\qual{Training and evaluation labels need appropriate follow-up.}
\talknote{0:50}{A conventional classifier can answer a perfectly legitimate fixed-horizon question, such as default within twelve months. The issue is not that classification is wrong: it is that a label requires the relevant horizon to be observed, or otherwise appropriately handled. An event observed before the horizon is known; a no-default label generally needs exposure through that horizon, with competing events addressed according to the intended target. Our six-month borrower does not yet supply that label. If already running long, skip the next worked-label slide and go directly to slide 7; this explanation is self-contained.}
\end{frame}

\begin{frame}[label=optional-labels]{Which 12-month labels are known?}
\optionalexplanation
\vspace{.5cm}
\begin{center}
\begin{tabular}{ll}\toprule
Observed history & Default by month 12?\\\midrule
Default at month 3 & \uncover<2->{Yes}\\[.3cm]
Default at month 24 & \uncover<2->{No}\\[.3cm]
No default through month 6 & \uncover<3->{Unknown}\\\bottomrule
\end{tabular}
\end{center}
\vspace{.3cm}\qual{Hypothetical loans; no earlier competing exit in these histories.}
\talknote{0:55}{OPTIONAL: skip this entire slide to save 55 seconds. It elaborates slide 5 without introducing a new concept. First read the three histories with the answers hidden. Reveal the first two labels: a default at month three is inside the twelve-month horizon, and a default at month twenty-four implies event-free exposure through month twelve in these simple histories. Pause, then reveal Unknown for the loan only observed through six months. This borrower is not a negative training example just because no default has yet appeared. Resume with slide 7.}
\end{frame}

\begin{frame}[label=why-survival]{Why survival analysis?}
\begin{itemize}\setlength\itemsep{.55cm}
\item<1-> Use when events occur.
\item<2-> Accommodate right-censored follow-up under suitable assumptions.
\item<3-> Here: Cox for interpretable relative default hazards.
\end{itemize}
\vspace{.35cm}\qual{Machine-learning methods can also model survival outcomes.}
\talknote{0:45}{Connect both opening problems to the method. Time-to-event data retain event timing and can use incomplete right-censored histories under appropriate observation assumptions. This does not mean censoring can be ignored or that survival analysis is automatically unbiased. We choose Cox because a small set of relative-hazard effects is interpretable for this lecture. Survival analysis is a type of outcome framework, not the opposite of machine learning; ML survival methods also exist. We will not compare their predictive performance.}
\end{frame}

\begin{frame}{Choose the event and the clock}
\centering
{\Large Origination $\longrightarrow$ first recorded default\par}
\vspace{.8cm}
Elapsed time since the loan was issued
\talknote{0:25}{Before writing any formula, specify the event and the start of the clock. Here the event is first recorded default and the origin is loan issuance. The illustrations use months; the audited example uses days. Some histories stop before a default is seen, and we need to distinguish why they stop. Move directly to the timelines.}
\end{frame}

\begin{frame}[plain]
\timelinefigure
\talknote{1:20}{Use the same hypothetical borrowers as in the opening: A defaults at three months and B at twenty-four. C is still active when observation ends at six months: this is right censoring. We know only that it remained default-free through the observed period. Add D, which closes at twelve months: closure is a competing event ending default exposure, not simply failure to keep observing an active loan. Point to the cross, censoring bar and closure diamond in turn. The whole figure is static; do not click expecting internal reveals. For interpreting the default rate, remove a closed loan from the exposed risk set. Appropriate administrative-censoring assumptions still matter.}
\end{frame}

\begin{frame}{The survival function}
\qual{First use a hypothetical default-only setting.}
\vspace{.6cm}
{\Large\[S(t)=P(T_D>t)\]}
\vspace{.6cm}
\uncover<2->{The probability of remaining default-free beyond time $t$.}
\talknote{0:40}{We now have a reason to introduce notation. Let T_D be time to default in a simplified world without competing closure. Read S(t) as the probability that the event has not happened by t. Reveal the plain-language interpretation after reading the equation. Survival is the technical term for remaining event-free; it does not imply anything about the borrower outside this defined event.}
\end{frame}

\begin{frame}[plain]
\fullfigure{02_illustrative_survival}
\talknote{1:00}{This curve is hypothetical, not a Bondora estimate. Start at one on the vertical axis and move right along elapsed months. Trace the month-twelve guide up to the curve, then across to about 0.79. Pause for three seconds so the audience can read the guides. About 79 percent remain default-free at twelve months in this default-only example. The formula exp(-0.02t) uses t in months. The same plot is already visible throughout this frame; guide attention verbally rather than pretending it animates.}
\end{frame}

\begin{frame}{How rapidly are defaults occurring now?}
Among loans still exposed now:\par
\vspace{.55cm}
how rapidly do defaults occur next?\par
\vspace{.7cm}
\uncover<2->{\textcolor{teal}{Hazard: an instantaneous rate, not a probability.}}
\talknote{0:35}{Introduce hazard in ordinary language before the limit definition. Restrict attention to loans still exposed at the current time and ask about the rate of imminent defaults. It is a local rate, not the fraction that will eventually default. Reveal the term hazard after that intuition. If time is measured in months, its units are inverse months. The next equation formalises this same question rather than introducing a different object.}
\end{frame}

\begin{frame}{From a short interval to a hazard rate}
\qual{Hypothetical default-only setting; $t$ in months.}
\vspace{.4cm}
\[
h(t)=\lim_{\Delta t\downarrow 0}
\frac{P(t\leq T_D<t+\Delta t\mid T_D\geq t)}{\Delta t}
\]
\vspace{.4cm}
\uncover<2->{\[
P(t\leq T_D<t+\Delta t\mid T_D\geq t)
\ \approx\ h(t)\,\Delta t
\]}
\uncover<2->{\textcolor{teal}{Rate $\times$ a short interval $\approx$ a conditional probability.}}
\talknote{1:05}{Read the numerator first: conditional on no default yet, what is the chance of default in the next small interval? Divide by its length and shrink the interval to obtain a rate. Do not explain every symbol as a separate reveal. Once the definition is explained, reveal the short-interval approximation. For the earlier curve a constant hazard of 0.02 per month gives an approximate 0.02 one-month conditional probability, with exact value 1-exp(-0.02), about 0.0198. Do not multiply a rate over a long horizon as though it were an exact probability. If running long, skip slide 14 and go directly to borrower characteristics on slide 15.}
\end{frame}

\begin{frame}[label=optional-cumulative]{From hazard to survival}
\optionalexplanation\quad\qual{Default-only setting.}
\vspace{.25cm}
{\Large\[H(t)=\int_0^t h(u)\,du\]}
\uncover<2->{{\Large\[S(t)=\exp\{-H(t)\}\]}}
\vspace{.35cm}
\uncover<2->{\qual{With competing closure, default hazard alone does not determine actual default probabilities.}}
\talknote{0:45}{OPTIONAL: skip this entire slide to save 45 seconds. It connects two already introduced objects. H accumulates the hazard over time and, in the default-only setting, S equals exp(-H). For the illustrative constant rate, H(t)=0.02t, giving the earlier curve. Reveal the survival equation only after the integral. With competing closure, exp(-H_D) is not the actual event-free survival or actual default probability: those depend on the competing process too. No cumulative-incidence analysis is being added. Continue with slide 15; it does not rely on H notation.}
\end{frame}

\begin{frame}[label=borrower-characteristics]{From timing to borrower characteristics}
We have described event timing.\par
\vspace{.8cm}
\uncover<2->{How does the default rate differ between borrowers?}
\talknote{0:25}{Make the progression explicit: We have described event timing. Now we want to understand how it differs between borrowers. Age, income or customer status may be associated with different rates. This motivates adding characteristics to a time-dependent baseline. Later the worked example will make the same distinction empirically: unadjusted timing curves describe accumulation, while Cox describes adjusted relative rates among exposed loans. No new regression is introduced.}
\end{frame}

\begin{frame}{The Cox model}
\vspace{.1cm}
{\Large\[
h(t\mid x)=\underbrace{h_0(t)}_{\text{baseline}}
\underbrace{\exp(\beta^{\mathsf T}x)}_{\text{multiplier}}
\]}
\vspace{.3cm}
\uncover<2->{Baseline: the time pattern for reference covariates.}\par
\vspace{.3cm}
\uncover<3->{Multiplier: how the covariates scale that rate.}\par
\vspace{.4cm}
\qual{For credit: a cause-specific default rate among loans still exposed.}
\talknote{1:20}{Show the whole equation first and read it as baseline times a positive covariate multiplier. Reveal the baseline explanation and point to h_0(t): we do not require a constant hazard or choose a fixed parametric shape here. It represents the rate at reference covariates. Then reveal the multiplier explanation: exp(beta transpose x) scales that rate. In the audited example the object is the cause-specific default hazard among loans neither defaulted nor closed, and the baseline is allowed to differ by origination year. Borrower coefficients are shared across those strata.}
\end{frame}

\begin{frame}{Compare otherwise identical borrowers}
\qual{Hypothetical new vs existing customer; same baseline and other covariates.}
\vspace{.45cm}
\[
\mathrm{HR}=\frac{h(t\mid\text{new})}{h(t\mid\text{existing})}
=\exp(\beta_{\mathrm{new}})
\]
\vspace{.5cm}
\uncover<2->{If HR $=1.5$: a 50\% higher rate at each $t$, under PH.}\par
\vspace{.4cm}
\qual{A conditional rate ratio, not a probability ratio or causal effect.}
\talknote{0:50}{The baseline cancels for two otherwise identical borrowers in the same baseline stratum. The customer indicator differs by one, so their relative hazard is exp(beta_new). Explain the equation before revealing the illustrative value of 1.5. That means a fifty percent higher instantaneous rate, not fifty percentage points or a fifty percent higher default probability over a horizon. This value is hypothetical and is not the final Bondora estimate. Under PH the ratio remains constant, even though both hazards may vary, as the next picture shows.}
\end{frame}

\begin{frame}[plain]
\fullfigure{03_proportional_hazards}
\talknote{0:50}{First follow A's curve, then B's. Both rates change over time; B stays one and a half times A. Point to the right panel showing the flat ratio. Pause briefly to let the audience compare the changing left panel with the constant right panel. Proportional hazards concerns the ratio, not constancy of either rate. These curves are hypothetical, not evidence that the assumption holds for Bondora. The plot is static throughout this frame.}
\end{frame}

\begin{frame}{How does Cox learn relative effects?}
\qual{Illustrative: three borrowers at risk just before one default.}
\vspace{.65cm}
\begin{columns}[T,onlytextwidth]
\column{.29\textwidth}\centering\textcolor{teal}{\Large A}\par\vspace{.3cm}
Existing\par customer\par\vspace{.4cm}Relative risk $1$
\column{.29\textwidth}\centering\textcolor{teal}{\Large B}\par\vspace{.3cm}
Existing\par customer\par\vspace{.4cm}Relative risk $1$
\column{.33\textwidth}\centering\textcolor{teal}{\Large C}\par\vspace{.3cm}
New\par customer\par\vspace{.4cm}Relative risk $e^\beta$
\end{columns}
\vspace{.6cm}
\centering Which borrower accounts for this default?
\talknote{1:00}{Introduce a separate teaching risk set, not the original three opening histories. All borrowers are still exposed just before this default time, in the same baseline stratum, and other covariates are identical. Their weights are one, one and exp(beta). Explain that we condition on one default occurring and ask which borrower it belongs to. Ask rhetorically what fraction of the total weight belongs to C, then pause for three seconds. Do not advance to the worked answer until the setup is understood.}
\end{frame}

\begin{frame}[plain]
\fullfigure{04_cox_risk_set}
\talknote{1:30}{Now show the answer: C's conditional share is exp(beta) divided by 2+exp(beta). Point to the symbolic expression before explaining the numerical example. With beta=log(2), chosen only for illustration, the weights are one, one and two, so C receives half of the total weight. Pause to allow the arithmetic to settle. This is conditional on one event occurring now, not a probability of default over a horizon. Cox combines comparisons like this across observed event times to estimate relative effects; the common baseline cancels. Full partial likelihood and the existing Efron handling of tied daily defaults are left to backup notes.}
\end{frame}

\begin{frame}{The audited Bondora cohort}
Estonia\quad$\cdot$\quad Issued 2013--2015\par
\vspace{.3cm}
Follow-up: at most 730 days from origination\par
\vspace{.6cm}
\uncover<2->{\begin{columns}[onlytextwidth]
\column{.32\textwidth}\centering{\Large\color{teal}\LoanCount}\par Loans
\column{.32\textwidth}\centering{\Large\color{teal}\BorrowerCount}\par Borrowers
\column{.32\textwidth}\centering{\Large\color{teal}\DefaultCount}\par Recorded defaults
\end{columns}}
\vspace{.5cm}
\qual{Reporting cutoff: 20 July 2021; all loans have 730 days of potential follow-up.}
\talknote{1:00}{Move clearly from hypothetical teaching examples to empirical data. The audited analysis retains 9,844 loans, 6,924 borrowers and 2,508 recorded defaults within 730 days. Reveal these counts after defining the country, origination years and horizon. The cutoff is the explicitly recorded ReportAsOfEOD of 20 July 2021, not the latest observed default or the file creation date. All origination dates have at least 730 days of potential observation. The event count is descriptive, not a model-predicted horizon probability.}
\end{frame}

\begin{frame}{Follow-up uses audited endpoints}
\begin{itemize}\setlength\itemsep{.45cm}
\item<1-> Default: first recorded \texttt{DefaultDate}.
\item<2-> Closure proxy: \texttt{ContractEndDate}\par
only when snapshot \texttt{Status=Repaid}.
\item<3-> Stop at the first exit or 730 days.
\end{itemize}
\vspace{.35cm}
\qual{Closure is not guaranteed full cash repayment. The source default threshold is unverified.}
\talknote{1:05}{Keep this operational definition short. A later repaid status does not erase an earlier recorded default. ContractEndDate occurs on Current and Late loans too and can be a future date, so not every such value is a repayment. The audit restricts the nondefault exit to a closure proxy in loans labeled Repaid; write-offs mean this still need not be literal full cash repayment. The general stopping rule also includes reporting cutoff, but for this cohort the 730-day cap comes first. Keep the uncertainty about the source's arrears threshold explicit; do not assert a sixty- or ninety-day definition.}
\end{frame}

\begin{frame}[plain]
\timingfigure
\talknote{1:50}{Move from endpoint definitions to when the recorded defaults accumulate. These are unadjusted descriptive cumulative-incidence curves, using Aalen–Johansen with the audited closure proxy as a competing event. The vertical axis is the estimated proportion of loans experiencing recorded default by each time before a competing closure. Point to twelve months: about 14.3 percent for new-customer loans and 9.9 percent for existing-customer loans. At twenty-four months the figures are about 27.4 and 21.2 percent. Read one horizon slowly rather than listing every number. The at-risk table includes observed loans immediately before each grid time. The month scale maps 730 days to 24 months. All have full potential follow-up, so these point estimates also agree with directly counted event fractions. No bands are shown: repeated borrowers are present and no independence-based uncertainty is being presented. New and existing customers differ in other covariates and vintages; these descriptive differences are not causal and are not the adjusted Cox HR. The next covariates/HR slides answer the complementary question about conditional rates.}
\end{frame}

\begin{frame}{Five predictors and their contrasts}
\begin{center}
\begin{tabular}{ll}
Age & $+10$ years\\[.28cm]
Income & Doubling\\[.28cm]
Amount & $+1{,}000$ recorded monetary units\\[.28cm]
Original scheduled term & $+12$ months\\[.28cm]
New-customer status & Yes vs no
\end{tabular}
\end{center}
\vspace{.3cm}
\qual{Year-stratified baseline; borrower-clustered intervals.}\par
\qual{Monthly original-schedule proxy; source currency/income period unverified.}
\talknote{0:55}{Read these as comparison units, not raw coefficient sizes. The model uses log income, and multiplying its coefficient by log(2) gives the displayed doubling contrast. Original scheduled term is reconstructed from FirstPaymentDate to MaturityDate_Original assuming monthly installments, not from revised maturity. Amount and income are used as recorded; the existing plot's euro label has the report's unverified-currency qualification. Baseline hazards differ by origination year and intervals cluster repeated borrowers. We use this one existing Cox model; no predictors, fits or selection procedures are added.}
\end{frame}

\begin{frame}[plain]
\resultfigure
\talknote{1:10}{Now ask which characteristics are associated with the conditional rate, adjusting for the included covariates and origination-year baseline. This cleaner figure reads the same final audited coefficients and borrower-clustered 95 percent intervals; none were refitted. Orient the audience to one and the logarithmic scale. Contrast it explicitly with the preceding unadjusted timing graph: these hazard ratios are adjusted relative rates among exposed loans, not differences or ratios of cumulative incidences. Exact predictor contrasts remain visible. Daily ties make the intervals approximate, and age/new-customer PH concerns remain visible. Detailed nonlinearity, diagnostic and source-definition caveats follow in limitations and backup notes.}
\end{frame}

\begin{frame}{New versus existing customer}
\centering
{\Large HR $=\NewCustomerHR$}\par
\vspace{.2cm}
95\% CI: $\NewCustomerLower$--$\NewCustomerUpper$\par
\vspace{.6cm}
\uncover<2->{About 44\% higher conditional default rate\par
among loans still exposed.}\par
\vspace{.45cm}
\uncover<2->{\qual{Association, not a default-probability ratio or causal effect.}\par
\qual{PH concern: this single HR is a model summary.}}
\talknote{1:10}{Show and explain the estimate first. On the next reveal, subtract one from the unrounded 1.4445 and convert to about forty-four percent. This is a conditional instantaneous default rate among loans still exposed, adjusting for the included covariates and origination-year baseline. It is not forty-four percentage points, a default-probability ratio or a causal effect of being new. The corrected diagnostics flag new-customer PH concerns in the full sample, so the single constant HR compresses a possibly time-varying association. Give that qualification before moving to general limitations.}
\end{frame}

\begin{frame}{Where interpretation is limited}
\begin{itemize}\setlength\itemsep{.55cm}
\item<1-> PH concerns: age and new-customer status.
\item<2-> Nonlinearity; approximate tied-date intervals.
\item<3-> Recorded endpoints and measurement proxies.
\end{itemize}
\talknote{1:20}{One short point per reveal. Corrected elapsed-rank and log-time diagnostics flag age and new-customer status; the superseded original-term PH finding is not repeated. Age also shows nonlinearity, while term has a borderline nonlinear departure, a separate assumption. Daily ties make the borrower-clustered Wald intervals approximate. The full-sample PH tests are approximate and not cluster adjusted. The source arrears threshold, full-repayment timing, monetary measurement definitions and original monthly-schedule proxy remain qualified. Detailed p-values and sensitivity checks are in backup material; do not turn this into a new modelling exercise.}
\end{frame}

\begin{frame}{When will a borrower default?}
\begin{itemize}\setlength\itemsep{.4cm}
\item<1-> No certain individual default date.
\item<2-> By 24 months: \NewCIFTwentyFour\% new vs \ExistingCIFTwentyFour\% existing.
\item<3-> Adjusted new-customer HR: $\NewCustomerHR$.
\end{itemize}
\vspace{.35cm}
\uncover<2->{\qual{Timing proportions are unadjusted/descriptive under the closure proxy.}}\par
\uncover<3->{\qual{HR is a conditional association, not a probability ratio or causal effect.}}\par
\vspace{.3cm}\uncover<3->{\textcolor{teal}{Questions}}
\talknote{1:00}{Return to the title and opening borrowers. A three-month default and a two-year default occur after different amounts of exposure and repayment; six observed default-free months do not reveal an eventual event date. We cannot assign certainty to an individual date, but we can describe accumulation: by 24 months about 27.4 percent of new-customer loans and 21.2 percent of existing-customer loans have a recorded default under the audited closure proxy. That comparison is unadjusted. Separately the unchanged Cox HR of 1.445 relates new-customer status to the adjusted conditional default rate, with its PH qualification. Cox is not inherently unable to estimate survival: with a fitted baseline it can do so in a default-only setting. Actual default probabilities under competing closure require accounting for both events. No predictive accuracy has been evaluated. End the speaking plan here and invite questions; questions and backups are outside the timing.}
\end{frame}
